\section{Memory区域与对象的存储期}

\begin{frame}{先猜输出：每次进入函数，变量都会重新开始吗？}
\begin{Code}*{c}\lstinputlisting{code/lifetime.c}\end{Code}
\note{先留一分钟，让同学写下两行输出；不要先把static等同于全局可见。}
\end{frame}

\begin{frame}{答案：作用域与生命周期是两件事}
\begin{itemize}
\item 第一次：\texttt{1 1 11}；第二次：\texttt{1 2 12}。
\item \texttt{local}：每次进入函数体，开始新的automatic对象生命周期。
\item \texttt{count}：具有static storage duration，整个程序执行期间都存在。
\item \texttt{global}：同样具有static storage duration，值不会在函数返回时重置。
\end{itemize}
\medskip\alert{count只能在函数体内用名字访问，但它不会随着函数返回而消失。}
\end{frame}

\begin{frame}{Scope、storage duration、lifetime}
\begin{itemize}
\item \textbf{Scope（作用域）}：源码中，哪些位置能用这个名字？
\item \textbf{Storage duration（存储期）}：对象的存储存在多长时间？
\item \textbf{Lifetime（生命周期）}：对象在执行中何时存在、何时可以访问？
\end{itemize}
\medskip
\begin{tabular}{lll}
\textbf{对象} & \textbf{典型存储期} & \textbf{何时结束}\\[4pt]
普通块内局部变量 & automatic & 离开对应的块\\
全局变量／\texttt{static}局部变量 & static & 程序结束\\
\texttt{malloc}分配的对象 & allocated & 被释放或被成功重新分配
\end{tabular}

\note{本节限定普通非VLA示例，不引入thread storage duration。指针也是对象，拥有自己的生命周期。}
\end{frame}

\begin{frame}{Memory layout：常见Linux进程的示意图}
\begin{columns}[c]
\begin{column}{0.59\textwidth}
\centering
\begin{tikzpicture}[x=1cm,y=1cm,>=Stealth,
    band/.style={draw=VSPPrimary!65,line width=0.55pt},
    name/.style={anchor=west,font=\small\ttfamily\bfseries,text=VSPPrimary},
    desc/.style={anchor=west,font=\small,text=black!80}]
    % One address space, with a fixed width for every region.
    \fill[VSPPrimary!7] (0,0) rectangle (6,-0.72);
    \fill[black!2] (0,-0.72) rectangle (6,-2.22);
    \fill[VSPPrimary!7] (0,-2.22) rectangle (6,-2.94);
    \fill[VSPPrimary!9] (0,-2.94) rectangle (6,-4.22);
    \fill[white] (0,-4.22) rectangle (6,-4.86);
    \draw[band] (0,0) rectangle (6,-4.86);
    \foreach \y in {-0.72,-2.22,-2.94,-3.58,-4.22}
        \draw[band] (0,\y) -- (6,\y);
    % Read-only objects and instructions have a neutral background.
    \fill[white] (0.01,-4.22) rectangle (5.99,-4.85);
    \draw[band] (0,-4.86) -- (6,-4.86);
    \node[name] at (0.24,-0.36) {Stack};
    \node[desc] at (1.85,-0.36) {调用栈、部分局部对象};
    \node[name] at (0.24,-2.58) {Heap};
    \node[desc] at (1.85,-2.58) {动态分配区域};
    \node[name] at (0.24,-3.26) {.bss};
    \node[desc] at (1.85,-3.26) {零初始化静态对象};
    \node[name] at (0.24,-3.90) {.data};
    \node[desc] at (1.85,-3.90) {非零初始化静态对象};
    \node[name] at (0.24,-4.54) {.rodata};
    \node[desc] at (1.85,-4.54) {字符串字面量等只读数据};
    % Group the usual writable static-storage sections without adding a section.
    \draw[VSPPrimary,line width=1pt] (0,-2.94) rectangle (6,-4.22);
    \draw[VSPPrimary,line width=0.9pt] (6.08,-2.94) -- (6.24,-2.94)
        -- (6.24,-4.22) -- (6.08,-4.22);
    \node[rotate=90,font=\scriptsize\bfseries,text=VSPPrimary] at (6.48,-3.58) {Static};
    % Code sits in its own equal-width region below the data regions.
    \fill[white] (0,-4.86) rectangle (6,-5.50);
    \draw[band] (0,-4.86) rectangle (6,-5.50);
    \node[name] at (0.24,-5.18) {.text};
    \node[desc] at (1.85,-5.18) {机器指令};
    % RAM/ROM are teaching analogies for writable objects vs read-only content.
    \draw[VSPPrimary,line width=1pt] (6.68,0) -- (6.84,0)
        -- (6.84,-4.22) -- (6.68,-4.22);
    \node[rotate=90,font=\scriptsize\bfseries,text=VSPPrimary] at (7.08,-2.11) {RAM（类比）};
    \draw[black!65,line width=1pt] (6.68,-4.22) -- (6.84,-4.22)
        -- (6.84,-5.50) -- (6.68,-5.50);
    \node[rotate=90,font=\scriptsize\bfseries,text=black!65] at (7.08,-4.86) {ROM（类比）};
    % The gap is not a segment; arrows connect growth to the right region.
    \draw[->,VSPPrimary,very thick] (0.55,-0.87) -- (0.55,-1.55);
    \node[anchor=west,font=\scriptsize,text=VSPPrimary] at (0.77,-1.12) {Stack向低地址增长};
    \draw[->,VSPPrimary,very thick] (5.45,-2.07) -- (5.45,-1.39);
    \node[anchor=east,font=\scriptsize,text=VSPPrimary] at (5.23,-1.83) {Heap向高地址增长};
    \node[font=\small,text=black!55,fill=black!2,inner sep=2pt] at (3,-1.47) {空闲空间／其他映射};
    % Address direction is independent from allocation growth.
    \draw[->,black!45,line width=0.7pt] (-0.42,-5.35) -- (-0.42,-0.15);
    \node[anchor=east,font=\scriptsize,text=black!65,align=center] at (-0.52,-0.30) {高\\地址};
    \node[anchor=east,font=\scriptsize,text=black!65,align=center] at (-0.52,-5.15) {低\\地址};
\end{tikzpicture}
\end{column}
\begin{column}{0.37\textwidth}
\textbf{再按读写用途分组}
\medskip
\begin{itemize}
    \item \textbf{RAM类比}：Stack、Heap及Static区中的可写对象。
    \item \textbf{ROM类比}：.rodata的只读数据与.text的机器指令。
    \item \textbf{Static区}：将.bss与.data归组，对象存储期贯穿程序。
\end{itemize}
\medskip
\small\alert{Linux中.text／.rodata通常也驻留在物理RAM中；图中ROM仅类比只读属性。}
\par\medskip
\small 地址由下向上增大；箭头表示常见增长方向。C标准不保证这些排列与方向。
\end{column}
\end{columns}
\note{RAM与ROM在本图中仅作按访问属性的教学类比，并非物理内存或硬件地址区间。实际Linux的.text通常可读可执行、.rodata通常只读，二者均可驻留物理RAM，权限由映射控制。空闲空间与其他映射不据此统一规定读写权限。Static是这里为教学使用的分组标签，不是另一个ELF section；.bss与.data分别保留。具有static存储期的只读对象也可能放在.rodata，故不把此分组说成所有static对象唯一所在的区域。}
\end{frame}

\begin{frame}{这些对象都在哪？没写初值时会怎样？}
\begin{columns}[T]
\begin{column}{0.45\textwidth}
\begin{Code}*{c}\lstinputlisting{code/locations.inc}\end{Code}
\end{column}
\begin{column}{0.51\textwidth}
\scriptsize
\setlength{\tabcolsep}{4pt}
\renewcommand{\arraystretch}{1.3}
\begin{tabular}{lll}
\textbf{对象／数据} & \textbf{初始值} & \textbf{通常所在区域}\\
\texttt{global\_uninit} & 0 & .bss（Static）\\
\texttt{global\_init} & 7 & .data（Static）\\
\texttt{static\_global} & 8 & .data（Static）\\
\texttt{static\_local} & 9 & .data（Static）\\
\texttt{static\_zero} & 0 & .bss（Static）\\
\texttt{local} & 不确定值 & Stack／寄存器\\
\texttt{p} & 字面量的地址 & Stack／寄存器\\
\texttt{"hello"} & 字符及结尾零 & .rodata
\end{tabular}
\end{column}
\end{columns}
\medskip
\begin{itemize}
    \item \textbf{默认初始化行为，看存储期：}没写初值的全局／static对象会\textbf{零初始化}；普通局部\texttt{int local;}不会自动置零，\alert{先赋值，再读取}。
    \item 非零初值通常进入.data；显式\texttt{= 0}也通常进入.bss。static局部对象的生命周期贯穿程序。
    \item \alert{p保存地址，"hello"是所指数据；字符串字面量不能修改。}
\end{itemize}
\note{这里的“默认初始化行为”是在C17语境中解释未显式初始化时会发生什么，不套用C++的default-initialization分类。普通automatic局部对象未初始化时具有不确定值，不以“随机数”或“垃圾值”描述，更不能读取后猜测结果。位置是常见实现，对象可能被优化；字符串字面量具有static存储期。}
\end{frame}

\begin{frame}{.bss与.data：为什么两个程序大小差这么多？}
\begin{columns}[T]\begin{column}{0.53\textwidth}\begin{Code}*{c}\lstinputlisting{code/bss.c}\end{Code}\begin{Code}*{c}\lstinputlisting{code/data.c}\end{Code}\end{column}\begin{column}{0.43\textwidth}\begin{itemize}
\item 两组数组在运行时都需要存储。
\item 第一组默认全为0；第二组前三项非零，其余仍为0。
\item .bss通常只记录需要多少零初始化空间。
\item .data需要在文件中保存初始数据。
\end{itemize}
\end{column}\end{columns}\medskip\alert{可执行文件小，不代表程序运行时不占内存。}
\end{frame}

\begin{frame}{现场看一眼：size与objdump}
在\texttt{r5/}目录执行：\begin{Code}*{bash}\lstinputlisting{code/layout.sh}\end{Code}\begin{itemize}
\item \texttt{size}：比较text、data、bss所占空间。
\item \texttt{objdump -t}：观察符号\texttt{array}属于哪个section。
\item \texttt{int a;}与\texttt{int b = 0;}通常都进入.bss。
\end{itemize}

\note{工具演示约3分钟，只观察section归属。零初始化不是未初始化；具体文件大小取决于平台与工具链。}
\end{frame}

\begin{frame}{Stack frame：函数返回后发生了什么？}

\begin{columns}[c]
\begin{column}{0.43\textwidth}
\begin{tikzpicture}[node distance=0.12cm,every node/.style={draw=VSPPrimary,minimum width=4.6cm,minimum height=0.85cm,align=center}]
\node (c) {\texttt{function\_c}的frame};
\node[below=of c] (b) {\texttt{function\_b}的frame};
\node[below=of b] {\texttt{main}的frame};
\end{tikzpicture}
\end{column}
\begin{column}{0.52\textwidth}
\begin{itemize}
\item 函数调用通常建立stack frame（栈帧）。
\item 栈帧可能保存局部变量、返回地址与部分参数。
\item 参数也可能通过寄存器传递。
\item 函数返回后，普通局部对象生命周期结束。
\end{itemize}
\end{column}
\end{columns}
\medskip\alert{内存里的旧字节暂时没被覆盖，也不代表原来的对象还活着。}

\end{frame}
